% Complete editable master: appendices/sound-history/student.md
\newcommand{\runningtitle}{Sound history}
% Layout inherited from the reviewed previous course; set a running title per unit.
% This is a small hand-maintained TeX layout, not a conversion pipeline.
\RequirePackage{fix-cm}
\documentclass[11pt,twoside]{article}
\usepackage{fontspec}
\setmainfont{Cambria}[
  Renderer=HarfBuzz,
  BoldFont=cambriab.ttf,
  ItalicFont=cambriai.ttf,
  BoldItalicFont=cambriaz.ttf
]
\setmonofont{lmmono10-regular.otf}[BoldFont=lmmonolt10-bold.otf]
\usepackage[paperwidth=165mm,paperheight=241mm,inner=21mm,outer=19mm,
  top=22mm,bottom=20mm,headheight=14pt,headsep=7mm,footskip=10mm]{geometry}
\usepackage{xcolor,array,booktabs,longtable,tabularx}
\usepackage{enumitem,fancyhdr}
\usepackage[unicode,hidelinks]{hyperref}
\definecolor{ink}{HTML}{000000}
\definecolor{muted}{HTML}{4C5558}
\setlength{\parindent}{0pt}
\setlength{\parskip}{3.5pt}
\AtBeginDocument{\fontsize{11}{14.3}\selectfont}
\setlength{\tabcolsep}{4pt}
\renewcommand{\arraystretch}{1.10}
\setlist[itemize]{leftmargin=1.3em,label=\textbullet,itemsep=2pt,topsep=3pt,parsep=0pt}
\setlist[enumerate]{leftmargin=1.8em,itemsep=3pt,topsep=4pt,parsep=0pt}
\widowpenalty=10000
\clubpenalty=10000
\emergencystretch=1em
\raggedbottom
\pagestyle{fancy}
\fancyhf{}
\providecommand{\runningtitle}{Old English}
\fancyhead[LE,RO]{\small\itshape \runningtitle}
\fancyhead[LO,RE]{\small Old English}
\fancyfoot[LE,RO]{\small\thepage}
\renewcommand{\headrulewidth}{0.3pt}
\newcommand{\booktitle}[2]{
  {\fontsize{9}{11}\selectfont\color{muted}#1\par}
  \vspace{3pt}
  {\fontsize{20}{24}\selectfont\color{ink}#2\par}
  \vspace{6pt}
}
\newcommand{\source}[1]{\par{\footnotesize\color{muted}#1\par}}
\newcommand{\readingtext}[1]{
  \par\addvspace{5pt}
  \begingroup\leftskip=1.5em
    \fontsize{12}{16}\selectfont #1\par
  \endgroup\addvspace{5pt}
}
\makeatletter
\renewcommand{\section}{\@startsection{section}{1}{\z@}
  {-2.4ex \@plus -0.6ex \@minus -0.2ex}{0.8ex \@plus 0.2ex}
  {\normalfont\fontsize{12}{15}\selectfont\bfseries\color{ink}}}
\renewcommand{\subsection}{\@startsection{subsection}{2}{\z@}
  {-1.8ex \@plus -0.5ex \@minus -0.2ex}{0.5ex \@plus 0.2ex}
  {\normalfont\fontsize{11}{14}\selectfont\bfseries\color{ink}}}
\makeatother

\hyphenpenalty=10000\exhyphenpenalty=10000
\emergencystretch=3em
\newcommand{\reserve}[1]{\par\penalty-100\begingroup\dimen0=\pagegoal\advance\dimen0 by-\pagetotal\ifdim\dimen0<#1\newpage\fi\endgroup}
\begin{document}
\addfontfeatures{Ligatures=TeXOff}
\raggedright
% appendix-block:B001:title:begin
\booktitle{APPENDIX}{Why Old English forms vary: a short sound history}
% appendix-block:B001:end
% appendix-block:B002:paragraph:begin
The course gives the forms needed to read and write Old English. This optional appendix explains some of their history. A change in a neighbouring sound, the loss of an ending, or the reshaping of a word after a familiar pattern can leave an alternation inside a paradigm. Breaking and umlaut are particularly useful explanations; several other processes help distinguish their results.\par
% appendix-block:B002:end
% appendix-block:B003:paragraph:begin
The examples retain the course’s spellings. Consult the pronunciation guide for the reading convention used in this course; the discussion here concerns how sounds changed over time.\par
% appendix-block:B003:end
% section: SH-S01
\reserve{5\baselineskip}
% appendix-block:B004:heading:begin
\section*{1. Reading a historical explanation}
% appendix-block:B004:end
% appendix-block:B005:paragraph:begin
A \emph{front vowel} is formed with the tongue relatively far forward in the mouth; a \emph{back vowel} has the tongue farther back. \emph{Rounding} means rounding the lips. A vowel can move forward without losing its lip rounding: this is relevant to the relation between \textbf{u} and \textbf{y}. \emph{Raising} means moving the tongue higher. These terms describe positions, not loudness or pitch.\par
% appendix-block:B005:end
% appendix-block:B006:paragraph:begin
A \emph{diphthong} changes position within one syllable. Old English distinguishes short and long diphthongs, just as it distinguishes short and long simple vowels. Two written vowel letters therefore do not by themselves indicate length. The macron marks the long forms: compare the spellings \textbf{ea} and \textbf{ēa}.\par
% appendix-block:B006:end
% appendix-block:B007:paragraph:begin
In historical explanations, \textbf{j} represents the consonant sound at the beginning of English \emph{yes}. A \emph{reconstructed form} is an earlier form inferred from evidence rather than preserved as an Old English example. Here a preceding asterisk marks a reconstruction. This local convention differs from Unit 3, where asterisks mark proposed errors in an exercise.\par
% appendix-block:B007:end
% appendix-block:B008:paragraph:begin
An arrow means “developed into”; it may summarize several changes rather than display every intermediate stage. Forms printed beside one another in a paradigm are coexisting grammatical forms. A singular does not become its plural through the historical passage of time.\par
% appendix-block:B008:end
% appendix-block:B009:paragraph:begin
Most of the explanations concern the prehistory of the West Saxon forms used in the course. Later changes and other dialects are labelled where they matter. A sound’s surroundings at the time of a change may differ from the letters visible in the surviving word.\par
% appendix-block:B009:end
% section: SH-S02
% examples: SH-E01
\reserve{5\baselineskip}
% appendix-block:B010:heading:begin
\section*{2. Fronting and restoration}
% appendix-block:B010:end
% appendix-block:B011:paragraph:begin
In stressed syllables, an early change brought short \textbf{a} forward to \textbf{æ}. This \emph{fronting}, also called \emph{Anglo-Frisian brightening}, did not apply in the same way before nasal consonants, \textbf{m} and \textbf{n}. The nasal environment had already given the vowel a different history.[1]\par
% appendix-block:B011:end
% appendix-block:B012:paragraph:begin
Later, some of the newly fronted \textbf{æ} returned to \textbf{a}. This \emph{restoration} occurred before an immediately following \textbf{w}, and when a back vowel in the next syllable followed an open first syllable. An \emph{open syllable} ends in a vowel. In the relevant pattern, a single consonant between vowels begins the following syllable; a cluster can close the first syllable and prevent restoration.\par
% appendix-block:B012:end
% appendix-block:B013:table:begin
\begingroup
\fontsize{10.5}{13.2}\selectfont
\setlength{\LTleft}{0pt}
\setlength{\LTright}{0pt plus 1fil}
\begin{longtable}{@{}>{\raggedright\arraybackslash}p{\dimexpr0.2500\linewidth-1.0000\tabcolsep\relax}>{\raggedright\arraybackslash}p{\dimexpr0.3500\linewidth-1.4000\tabcolsep\relax}>{\raggedright\arraybackslash}p{\dimexpr0.4000\linewidth-1.6000\tabcolsep\relax}@{}}
\toprule
\textbf{Form} & \textbf{Meaning and identity} & \textbf{What the form illustrates} \\
\midrule
\endfirsthead
\toprule
\textbf{Form} & \textbf{Meaning and identity} & \textbf{What the form illustrates} \\
\midrule
\endhead
\bottomrule
\endfoot
\bottomrule
\endlastfoot
\textbf{dæg} & day; nominative singular & The front vowel remains. \\
\textbf{dagas} & days; nominative plural & A following back vowel allowed restoration in the open first syllable. \\
\textbf{cræftum} & with skills; dative plural & The consonant cluster closes the first syllable, so the back vowel of the ending does not produce the same result. \\
\end{longtable}
\endgroup
% appendix-block:B013:end
% appendix-block:B014:paragraph:begin
Thus the familiar singular/plural alternation has a historical explanation. Its present distribution is learned with the noun: changing every \textbf{æ} to \textbf{a} before an ending would erase the distinction illustrated by \textbf{cræftum}.\par
% appendix-block:B014:end
% section: SH-S03
% examples: SH-E02
\reserve{5\baselineskip}
% appendix-block:B015:heading:begin
\section*{3. Breaking}
% appendix-block:B015:end
% appendix-block:B016:paragraph:begin
\emph{Breaking} is the development of a diphthong from a stressed front vowel before particular consonants. A second vowel element arose as the pronunciation moved toward the following consonant. The immediate short-vowel inputs were \textbf{æ}, \textbf{e} and \textbf{i}. Earlier \textbf{a} had already been fronted to \textbf{æ} where this was relevant.[2]\par
% appendix-block:B016:end
% appendix-block:B017:paragraph:begin
The conditions differ according to the vowel. In the following table, “another consonant” includes a second \textbf{r}, but excludes \textbf{j} after \textbf{r}.\par
% appendix-block:B017:end
% appendix-block:B018:table:begin
\begingroup
\fontsize{10.5}{13.2}\selectfont
\setlength{\LTleft}{0pt}
\setlength{\LTright}{0pt plus 1fil}
\begin{longtable}{@{}>{\raggedright\arraybackslash}p{\dimexpr0.2500\linewidth-1.0000\tabcolsep\relax}>{\raggedright\arraybackslash}p{\dimexpr0.3500\linewidth-1.4000\tabcolsep\relax}>{\raggedright\arraybackslash}p{\dimexpr0.4000\linewidth-1.6000\tabcolsep\relax}@{}}
\toprule
\textbf{Earlier short stressed vowel} & \textbf{Result} & \textbf{Principal environments} \\
\midrule
\endfirsthead
\toprule
\textbf{Earlier short stressed vowel} & \textbf{Result} & \textbf{Principal environments} \\
\midrule
\endhead
\bottomrule
\endfoot
\bottomrule
\endlastfoot
\textbf{æ} & \textbf{ea} & Before \textbf{h}; before \textbf{r} plus another consonant; before \textbf{l} plus another consonant. \\
\textbf{e} & \textbf{eo} & Before \textbf{h}; before \textbf{r} plus another consonant. Before \textbf{l} plus a consonant, the regular case is \textbf{lh}. \\
\textbf{i} & \textbf{io} & Before \textbf{h} or \textbf{r} plus another consonant. Breaking before \textbf{l} is not securely established. \\
\end{longtable}
\endgroup
% appendix-block:B018:end
% appendix-block:B019:paragraph:begin
At the time of breaking, \textbf{h} represented a friction sound made farther back in the mouth, rather than the weak initial sound familiar from English \emph{house}. The exclusion of \textbf{rj} and the restricted treatment before \textbf{l} are part of the rule. “A front vowel before a consonant cluster” is too broad a description.\par
% appendix-block:B019:end
% appendix-block:B020:paragraph:begin
The contrast between \textbf{helpan}, “help”, and \textbf{healp}, “helped” (preterite singular), makes the vowel-specific condition visible. The reconstructed earlier preterite had \textbf{æ}:\par
% appendix-block:B020:end
% appendix-block:B021:paragraph:begin
\textbf{*hælp} → \textbf{healp}, “helped”.\par
% appendix-block:B021:end
% appendix-block:B022:paragraph:begin
That vowel broke before \textbf{lp}. The \textbf{e} in \textbf{helpan} did not break before the same cluster: regular breaking of \textbf{e} before \textbf{l} requires the more restricted environment just given. The two forms began with different stem vowels; breaking affected them differently. The adjective \textbf{eald}, “old”, likewise shows the outcome of \textbf{æ} before \textbf{ld}.\par
% appendix-block:B022:end
% appendix-block:B023:paragraph:begin
Breaking a short vowel did not make it long. The \textbf{ea} in \textbf{healp} and \textbf{eald} is short. Their following consonant clusters are a separate matter, relevant to syllable weight (§7). Some long front vowels also underwent breaking before \textbf{h}; the table above is a summary of short-vowel breaking, not its whole history.\par
% appendix-block:B023:end
% appendix-block:B024:paragraph:begin
West Saxon subsequently merged older \textbf{io} with \textbf{eo}. The same written \textbf{eo} can consequently have more than one earlier vowel behind it. Nor does every \textbf{ea} or \textbf{eo} come from breaking: §5 describes other sources of diphthong spellings.\par
% appendix-block:B024:end
% section: SH-S04
\reserve{5\baselineskip}
% appendix-block:B025:heading:begin
\section*{4. Umlaut: the influence of a following syllable}
% appendix-block:B025:end
% appendix-block:B026:paragraph:begin
\emph{Umlaut}, or \emph{mutation}, is a change in one vowel under the influence of another sound in a following syllable. The most useful kind here is \emph{i-umlaut}, also called \emph{front mutation}. An older \textbf{i} or \textbf{j}, usually in the next syllable, caused the stressed vowel to be pronounced farther forward or higher. The speaker anticipated part of the following sound’s tongue position.[3]\par
% appendix-block:B026:end
% appendix-block:B027:paragraph:begin
The conditioning \textbf{i} or \textbf{j} often disappeared later, or was replaced by another vowel. Its effect on the root could remain. An Old English form may therefore have a mutated vowel even though its surviving ending contains neither \textbf{i} nor \textbf{j}.\par
% appendix-block:B027:end
% appendix-block:B028:paragraph:begin
Some characteristic results are:\par
% appendix-block:B028:end
% appendix-block:B029:table:begin
\begingroup
\fontsize{10.5}{13.2}\selectfont
\setlength{\LTleft}{0pt}
\setlength{\LTright}{0pt plus 1fil}
\begin{longtable}{@{}>{\raggedright\arraybackslash}p{\dimexpr0.2500\linewidth-1.0000\tabcolsep\relax}>{\raggedright\arraybackslash}p{\dimexpr0.3500\linewidth-1.4000\tabcolsep\relax}>{\raggedright\arraybackslash}p{\dimexpr0.4000\linewidth-1.6000\tabcolsep\relax}@{}}
\toprule
\textbf{Earlier vowel} & \textbf{Front-mutation result} & \textbf{Qualification} \\
\midrule
\endfirsthead
\toprule
\textbf{Earlier vowel} & \textbf{Front-mutation result} & \textbf{Qualification} \\
\midrule
\endhead
\bottomrule
\endfoot
\bottomrule
\endlastfoot
\textbf{u}, \textbf{ū} & \textbf{y}, \textbf{ȳ} & The tongue moves forward while the lips remain rounded. \\
\textbf{o}, \textbf{ō} & \textbf{e}, \textbf{ē} & Front vowels replace the back rounded vowels. \\
\textbf{æ} & \textbf{e} & The already front vowel is raised. \\
\textbf{ā} & \textbf{ǣ} & This compact correspondence excludes the separate nasal-vowel history. \\
\textbf{ea}, \textbf{ēa} & \textbf{ie}, \textbf{īe} & West Saxon diphthong correspondences; later spellings can differ. \\
\end{longtable}
\endgroup
% appendix-block:B029:end
% appendix-block:B030:paragraph:begin
This is a historical guide to selected correspondences. It is not an instruction to replace a vowel whenever a later written ending happens to contain \textbf{i}. The relevant environment must have existed when the change operated.\par
% appendix-block:B030:end
% examples: SH-E03
\reserve{5\baselineskip}
% appendix-block:B031:subheading:begin
\subsection*{Nouns}
% appendix-block:B031:end
% appendix-block:B032:paragraph:begin
\textbf{Fōt}, “foot”, has \textbf{fēt}, “feet”, in the nominative and accusative plural. The same \textbf{fēt} is also the selected dative singular, where it still means “foot”. The long \textbf{ē} reflects a formerly front-mutation environment. Other forms of the noun retain \textbf{ō}.\par
% appendix-block:B032:end
% appendix-block:B033:paragraph:begin
This explains an alternation already learned in the paradigm. Umlaut is not a general Old English plural ending: it occurs in particular inherited formations, and later analogy can also affect its distribution.\par
% appendix-block:B033:end
% examples: SH-E04
\reserve{5\baselineskip}
% appendix-block:B034:subheading:begin
\subsection*{Related words}
% appendix-block:B034:end
% appendix-block:B035:paragraph:begin
The noun \textbf{dōm}, “judgment”, and the weak verb \textbf{dēman}, “judge”, belong to the same word family. The earlier formation of the verb included \textbf{j}. Fulk represents that earlier verb as \textbf{*dōm(i)jan}, “judge”:\par
% appendix-block:B035:end
% appendix-block:B036:paragraph:begin
\textbf{*dōm(i)jan} → \textbf{dēman}, “judge”.\par
% appendix-block:B036:end
% appendix-block:B037:paragraph:begin
The arrow summarizes front mutation and the subsequent changes that removed the conditioning sound. The noun and the verb did not have identical endings, so their root vowels need not have developed identically. The surviving \textbf{-an} of \textbf{dēman} conceals the earlier conditioning sound.\par
% appendix-block:B037:end
% examples: SH-E05
\reserve{5\baselineskip}
% appendix-block:B038:subheading:begin
\subsection*{Comparison}
% appendix-block:B038:end
% appendix-block:B039:table:begin
\begingroup
\fontsize{10.5}{13.2}\selectfont
\setlength{\LTleft}{0pt}
\setlength{\LTright}{0pt plus 1fil}
\begin{longtable}{@{}>{\raggedright\arraybackslash}p{\dimexpr0.2500\linewidth-1.0000\tabcolsep\relax}>{\raggedright\arraybackslash}p{\dimexpr0.3500\linewidth-1.4000\tabcolsep\relax}>{\raggedright\arraybackslash}p{\dimexpr0.4000\linewidth-1.6000\tabcolsep\relax}@{}}
\toprule
\textbf{Positive} & \textbf{Comparative} & \textbf{Superlative} \\
\midrule
\endfirsthead
\toprule
\textbf{Positive} & \textbf{Comparative} & \textbf{Superlative} \\
\midrule
\endhead
\bottomrule
\endfoot
\bottomrule
\endlastfoot
\textbf{eald}, old & \textbf{ieldra}, older & \textbf{ieldest}, oldest \\
\end{longtable}
\endgroup
% appendix-block:B039:end
% appendix-block:B040:paragraph:begin
The comparison forms had an older front-mutation environment absent from the positive. Their \textbf{ie} belongs with the \textbf{ea} of \textbf{eald}. This family also shows that processes can operate in succession: breaking supplied one diphthong, and later mutation could change it. The comparative shown is masculine nominative singular; the table compares stems rather than replacing the complete adjective paradigms.\par
% appendix-block:B040:end
% examples: SH-E06
\reserve{5\baselineskip}
% appendix-block:B041:subheading:begin
\subsection*{Present singulars, and an older change}
% appendix-block:B041:end
% appendix-block:B042:paragraph:begin
Older second- and third-person singular present endings contained \textbf{i}. This helps explain why a changed vowel can occur in those cells while the infinitive, first singular and plural retain another vowel. Compare \textbf{hātan}, “call; name”, with third singular \textbf{hǣtt}, “calls; names”. Here \textbf{ā} beside \textbf{ǣ} illustrates front mutation; the consonant ending has its own development.\par
% appendix-block:B042:end
% appendix-block:B043:paragraph:begin
The course also places \textbf{helpan}, “help”, beside \textbf{hilpð}, “helps”. Its historical explanation needs a distinction. The raising of \textbf{e} to \textbf{i} before an older \textbf{i} or \textbf{j} took place earlier than the later front-mutation changes just described. Cook and Fulk both recognize this distinction. A practical description of “present mutation” can group the visible alternations together without making them one event in time.[4]\par
% appendix-block:B043:end
% appendix-block:B044:paragraph:begin
The familiar \textbf{ēo} beside \textbf{īe} in some present stems likewise involves an earlier raised vowel and further development. A one-step equation between the written Old English vowels leaves out that history. Keep the learned form in its actual person, tense and mood; a historical explanation does not extend an indicative singular alternation to an entire subjunctive paradigm.\par
% appendix-block:B044:end
% section: SH-S05
\reserve{5\baselineskip}
% appendix-block:B045:heading:begin
\section*{5. Other origins and later treatment of diphthongs}
% appendix-block:B045:end
% examples: SH-E07
\reserve{5\baselineskip}
% appendix-block:B046:subheading:begin
\subsection*{After an initial palatal consonant}
% appendix-block:B046:end
% appendix-block:B047:paragraph:begin
A \emph{palatal consonant} is pronounced with the tongue toward the hard palate, the front part of the roof of the mouth. In West Saxon, an initial palatal consonant could affect the stressed vowel following it. This \emph{palatal diphthongization} changed \textbf{æ}, \textbf{ǣ} to \textbf{ea}, \textbf{ēa}, and \textbf{e}, \textbf{ē} to \textbf{ie}, \textbf{īe}, in the relevant environments.[5]\par
% appendix-block:B047:end
% appendix-block:B048:paragraph:begin
The direction of influence differs from breaking: the vowel follows the conditioning consonant. The change also differs from the consonant itself becoming palatal, discussed in §6. It followed breaking; a vowel already altered by breaking did not simply undergo the same conversion again.\par
% appendix-block:B048:end
% appendix-block:B049:paragraph:begin
\textbf{Geat}, “gate”, illustrates the vowel effect after an initial palatal. Its plural \textbf{gatu}, “gates”, retains a back vowel. The earlier distinction between front and restored back vowels (§2) is therefore part of the explanation; the written letter \textbf{g} alone cannot predict the vowel.\par
% appendix-block:B049:end
% appendix-block:B050:paragraph:begin
The familiar verb provides several related results:\par
% appendix-block:B050:end
% appendix-block:B051:table:begin
\begingroup
\fontsize{10.5}{13.2}\selectfont
\setlength{\LTleft}{0pt}
\setlength{\LTright}{0pt plus 1fil}
\begin{longtable}{@{}>{\raggedright\arraybackslash}p{\dimexpr0.2500\linewidth-1.5000\tabcolsep\relax}>{\raggedright\arraybackslash}p{\dimexpr0.2500\linewidth-1.5000\tabcolsep\relax}>{\raggedright\arraybackslash}p{\dimexpr0.2500\linewidth-1.5000\tabcolsep\relax}>{\raggedright\arraybackslash}p{\dimexpr0.2500\linewidth-1.5000\tabcolsep\relax}@{}}
\toprule
\textbf{Infinitive} & \textbf{Preterite singular} & \textbf{Preterite plural} & \textbf{Past participle} \\
\midrule
\endfirsthead
\toprule
\textbf{Infinitive} & \textbf{Preterite singular} & \textbf{Preterite plural} & \textbf{Past participle} \\
\midrule
\endhead
\bottomrule
\endfoot
\bottomrule
\endlastfoot
\textbf{giefan}, give & \textbf{geaf}, gave & \textbf{gēafon}, gave & \textbf{giefen}, given \\
\end{longtable}
\endgroup
% appendix-block:B051:end
% appendix-block:B052:paragraph:begin
These are four principal parts used in the same language, not successive stages of a change. In particular, \textbf{geaf} has short \textbf{ea}, whereas \textbf{gēafon} has long \textbf{ēa}. Palatal diphthongization preserves that distinction of length. The pattern described here is centred on West Saxon; other Old English dialects need not show the same results.\par
% appendix-block:B052:end
% examples: SH-E08
\reserve{5\baselineskip}
% appendix-block:B053:subheading:begin
\subsection*{Back mutation}
% appendix-block:B053:end
% appendix-block:B054:paragraph:begin
\emph{Back mutation} is another influence from a following syllable. A back vowel could cause a short stressed front vowel to acquire a second element: earlier \textbf{i} and \textbf{e} could give \textbf{io} and \textbf{eo}. With the West Saxon merger of \textbf{io} and \textbf{eo}, the usual visible result is \textbf{eo}. \textbf{Heofon}, “sky; heaven”, is an example.[6]\par
% appendix-block:B054:end
% appendix-block:B055:paragraph:begin
West Saxon restricts the change chiefly to a single intervening \emph{labial}, a consonant made with the lips, or a \emph{liquid}, \textbf{l} or \textbf{r}. The relevant labials are \textbf{p}, \textbf{f}, \textbf{m} and \textbf{w}. More than one intervening consonant normally prevents this West Saxon change. The following vowels are also unequal in effect: \textbf{a} affects preceding \textbf{i}, but does not produce this mutation of preceding \textbf{e} in West Saxon.\par
% appendix-block:B055:end
% appendix-block:B056:paragraph:begin
The distribution is wider in other dialects. Even within West Saxon, speakers could restore an unmutated vowel by following related forms. Such reshaping is called \emph{analogy}. The source spelling \textbf{hefon}, also “sky; heaven”, has a simple vowel; for dictionary lookup, see \textbf{heofon}. A historical tendency can thus coexist with forms that have been altered after it operated.\par
% appendix-block:B056:end
% examples: SH-E15
\reserve{5\baselineskip}
% appendix-block:B057:subheading:begin
\subsection*{Smoothing}
% appendix-block:B057:end
% appendix-block:B058:paragraph:begin
\emph{Smoothing} replaces a diphthong with a single vowel. Late West Saxon smoothing affects \textbf{ea}, \textbf{ēa} in certain environments, including before \textbf{c}, \textbf{g}, \textbf{h}, or after an initial palatal. It is inconsistently represented in writing. A labelled example is \textbf{hēah}, “high”, beside the Late West Saxon spelling \textbf{hēh}, also “high”. For \textbf{hēh}, see \textbf{hēah} in the glossary.[7]\par
% appendix-block:B058:end
% appendix-block:B059:paragraph:begin
\emph{Anglian smoothing} has a different and wider scope. It affects more diphthongs before back consonants, sometimes with \textbf{r} or \textbf{l} intervening. The two developments should retain their dialect and period labels even where a resulting spelling looks alike. Neither supplies a universal instruction to simplify Old English diphthongs.\par
% appendix-block:B059:end
% section: SH-S06
\reserve{5\baselineskip}
% appendix-block:B060:heading:begin
\section*{6. Consonants and the order of changes}
% appendix-block:B060:end
% examples: SH-E09
\reserve{5\baselineskip}
% appendix-block:B061:subheading:begin
\subsection*{Palatalization}
% appendix-block:B061:end
% appendix-block:B062:paragraph:begin
\emph{Palatalization} changes a consonant’s articulation toward the hard palate. In appropriate historical environments, consonants written \textbf{c} or \textbf{g} acquired the sounds represented by English \textbf{ch} or consonantal \textbf{y}. The pronunciation guide gives their practical values and exceptions. Their history explains why a following front-vowel letter is useful evidence without being a complete spelling rule.[8]\par
% appendix-block:B062:end
% appendix-block:B063:paragraph:begin
Consider \textbf{cēlan}, “cool”, whose \textbf{c} is pronounced like English \textbf{k}. Its earlier form \textbf{*kōljan}, “cool”, had a back vowel next to the initial consonant. The later fronting produced \textbf{ē} after the consonant had remained nonpalatal. Reading the surviving \textbf{ē} as proof that the consonant must have become palatal reverses the order of events.\par
% appendix-block:B063:end
% appendix-block:B064:paragraph:begin
Different consonants and positions had different conditions, and some details remain disputed. The useful principle is chronological: a sound change responded to the sounds present at its own date, including sounds later lost or altered.\par
% appendix-block:B064:end
% examples: SH-E10
\reserve{5\baselineskip}
% appendix-block:B065:subheading:begin
\subsection*{Gemination}
% appendix-block:B065:end
% appendix-block:B066:paragraph:begin
\emph{Gemination} is consonant lengthening. A following \textbf{j} commonly lengthened the preceding consonant after a light syllable (see §7); \textbf{r} is an important exception. The \textbf{j} could later disappear while its effect remained.[9]\par
% appendix-block:B066:end
% appendix-block:B067:paragraph:begin
Old English \textbf{biddan}, “ask”, has a doubled consonant. Compare Gothic \textbf{bidjan}, also “ask”, which visibly retains \textbf{j}. Gothic is a related Germanic language: its form is a \emph{cognate}, not an Old English spelling or the parent of the Old English word. The comparison makes the historical conditioning sound easier to recognize. Other processes can also produce doubled consonants, so a double letter alone does not prove an earlier \textbf{j}.\par
% appendix-block:B067:end
% section: SH-S07
\reserve{5\baselineskip}
% appendix-block:B068:heading:begin
\section*{7. Vowel loss, weight and contraction}
% appendix-block:B068:end
% examples: SH-E11
\reserve{5\baselineskip}
% appendix-block:B069:subheading:begin
\subsection*{Final and internal vowels}
% appendix-block:B069:end
% appendix-block:B070:paragraph:begin
\emph{Apocope} is loss at the end of a word; \emph{syncope} is loss inside it. Old English lost many unstressed vowels, but the results depended on their position and on syllable weight. Vowel length and weight are distinct: a short vowel can occur in a heavy stem.[10]\par
% appendix-block:B070:end
% appendix-block:B071:paragraph:begin
For the ordinary noun pattern illustrated here, a short vowel or short diphthong followed by one consonant gives a light stem. A long vowel or long diphthong gives a heavy stem; so does a short vowel followed by two consonants. A doubled consonant counts as two. These differences helped determine whether a final short high vowel survived.\par
% appendix-block:B071:end
% appendix-block:B072:table:begin
\begingroup
\fontsize{10.5}{13.2}\selectfont
\setlength{\LTleft}{0pt}
\setlength{\LTright}{0pt plus 1fil}
\begin{longtable}{@{}>{\raggedright\arraybackslash}p{\dimexpr0.3500\linewidth-1.4000\tabcolsep\relax}>{\raggedright\arraybackslash}p{\dimexpr0.3500\linewidth-1.4000\tabcolsep\relax}>{\raggedright\arraybackslash}p{\dimexpr0.3000\linewidth-1.2000\tabcolsep\relax}@{}}
\toprule
\textbf{Nominative/accusative singular} & \textbf{Nominative/accusative plural} & \textbf{Relevant contrast} \\
\midrule
\endfirsthead
\toprule
\textbf{Nominative/accusative singular} & \textbf{Nominative/accusative plural} & \textbf{Relevant contrast} \\
\midrule
\endhead
\bottomrule
\endfoot
\bottomrule
\endlastfoot
\textbf{hof}, dwelling & \textbf{hofu}, dwellings & Light stem: plural \textbf{-u} remains. \\
\textbf{word}, word & \textbf{word}, words & Short \textbf{o}, but a heavy stem because of \textbf{rd}. \\
\textbf{wīf}, wife & \textbf{wīf}, wives & Heavy stem with long \textbf{ī}. \\
\end{longtable}
\endgroup
% appendix-block:B072:end
% appendix-block:B073:paragraph:begin
Historical high-vowel loss helps explain this distribution; it does not remove all final vowels from every heavy word. Later analogy and the histories of individual noun classes also matter. Use the actual paradigm for a particular noun.\par
% appendix-block:B073:end
% appendix-block:B074:paragraph:begin
For internal loss, compare \textbf{hēafod}, “head”, with genitive \textbf{hēafdes}, “head’s”. The unstressed middle vowel is absent before the ending. The source also records full \textbf{hēafodes}, “head’s”, where analogy has restored the fuller stem. For \textbf{hēafodes}, see \textbf{hēafod} in the glossary. The coexistence of the two genitives shows why a broad syncope rule cannot replace lexical and period-specific evidence.\par
% appendix-block:B074:end
% examples: SH-E12
\reserve{5\baselineskip}
% appendix-block:B075:subheading:begin
\subsection*{Loss of a consonant and contraction}
% appendix-block:B075:end
% appendix-block:B076:paragraph:begin
When an intervocalic \textbf{h} disappeared, it brought vowels together. \emph{Contraction} could then remove the second vowel and leave a long vowel or diphthong. Thus the reconstructed earlier infinitive \textbf{*seohan}, “see”, developed into \textbf{sēon}, “see”.[11]\par
% appendix-block:B076:end
% appendix-block:B077:paragraph:begin
The imperative singular \textbf{seoh}, “see!”, retains \textbf{h} in a different environment. Its \textbf{eo} is short; the infinitive’s \textbf{ēo} is long. Contraction therefore explains an important part of the pattern without accounting by itself for every form of a contracted verb. Loss of a consonant can also cause \emph{compensatory lengthening}, an increase in the preceding vowel’s duration; the exact environment must be checked separately.\par
% appendix-block:B077:end
% section: SH-S08
\reserve{5\baselineskip}
% appendix-block:B078:heading:begin
\section*{8. Ablaut, inherited alternations and analogy}
% appendix-block:B078:end
% examples: SH-E13
% appendix-block:B079:paragraph:begin
\emph{Ablaut} is an inherited alternation of stem vowels. Strong-verb principal parts preserve such alternations, subsequently reshaped by other sound changes. They are not formed by applying umlaut to an infinitive when a past meaning is needed.[12]\par
% appendix-block:B079:end
% appendix-block:B080:table:begin
\begingroup
\fontsize{10.5}{13.2}\selectfont
\setlength{\LTleft}{0pt}
\setlength{\LTright}{0pt plus 1fil}
\begin{longtable}{@{}>{\raggedright\arraybackslash}p{\dimexpr0.2500\linewidth-1.5000\tabcolsep\relax}>{\raggedright\arraybackslash}p{\dimexpr0.2500\linewidth-1.5000\tabcolsep\relax}>{\raggedright\arraybackslash}p{\dimexpr0.2500\linewidth-1.5000\tabcolsep\relax}>{\raggedright\arraybackslash}p{\dimexpr0.2500\linewidth-1.5000\tabcolsep\relax}@{}}
\toprule
\textbf{Infinitive} & \textbf{Preterite singular} & \textbf{Preterite plural} & \textbf{Past participle} \\
\midrule
\endfirsthead
\toprule
\textbf{Infinitive} & \textbf{Preterite singular} & \textbf{Preterite plural} & \textbf{Past participle} \\
\midrule
\endhead
\bottomrule
\endfoot
\bottomrule
\endlastfoot
\textbf{drīfan}, drive & \textbf{drāf}, drove & \textbf{drifon}, drove & \textbf{drifen}, driven \\
\textbf{helpan}, help & \textbf{healp}, helped & \textbf{hulpon}, helped & \textbf{holpen}, helped \\
\end{longtable}
\endgroup
% appendix-block:B080:end
% appendix-block:B081:paragraph:begin
The first series has \textbf{ī}, \textbf{ā}, \textbf{i}, \textbf{i}. The second shows the combined effects of inherited alternation and later change: the earlier different vowel of the singular preterite underwent breaking, giving \textbf{healp} (§3). Its \textbf{ea} is therefore not an untouched original ablaut vowel. The present singular \textbf{hilpð}, “helps”, belongs to the separate raising history explained in §4.\par
% appendix-block:B081:end
% appendix-block:B082:paragraph:begin
One word family can preserve several layers of history. A different strong class need not have the same series, and a weak verb can also have a changed stem vowel while retaining its dental preterite. The principal parts identify the forms actually needed.\par
% appendix-block:B082:end
% examples: SH-E14
% appendix-block:B083:paragraph:begin
Some inherited alternations involve consonants. Compare \textbf{cweðan}, “say”, and its past participle \textbf{cweden}, “said”. The \textbf{ð} beside \textbf{d} belongs to a history of consonant alternation. \emph{Verner’s law} describes an early accent-conditioned voicing process whose results were further altered by later sound changes. It operated before the ordinary Germanic stress pattern became established.[13]\par
% appendix-block:B083:end
% appendix-block:B084:paragraph:begin
This inherited alternation is distinct from the contextual voicing of fricatives in ordinary Old English pronunciation. It neither makes \textbf{ð} and \textbf{d} interchangeable spellings nor allows the present position of stress to generate all the older forms. Analogy can preserve, extend or remove an alternation, just as it can restore a lost vowel. Historical explanation and the recorded paradigm are therefore complementary kinds of evidence.\par
% appendix-block:B084:end
% section: SH-S09
\reserve{5\baselineskip}
% appendix-block:B085:heading:begin
\section*{9. Reference summary and source notes}
% appendix-block:B085:end
% appendix-block:B086:table:begin
\begingroup
\fontsize{10.5}{13.2}\selectfont
\setlength{\LTleft}{0pt}
\setlength{\LTright}{0pt plus 1fil}
\begin{longtable}{@{}>{\raggedright\arraybackslash}p{\dimexpr0.3400\linewidth-0.6800\tabcolsep\relax}>{\raggedright\arraybackslash}p{\dimexpr0.6600\linewidth-1.3200\tabcolsep\relax}@{}}
\toprule
\textbf{Process} & \textbf{Historical relation to look for} \\
\midrule
\endfirsthead
\toprule
\textbf{Process} & \textbf{Historical relation to look for} \\
\midrule
\endhead
\bottomrule
\endfoot
\bottomrule
\endlastfoot
Fronting and restoration & A stressed front vowel and a restored back vowel, conditioned by an earlier following sound and syllable structure. \\
Breaking & A stressed front vowel acquiring a second element before particular consonants; the conditions differ by vowel. \\
Front mutation & A vowel anticipating older \textbf{i} or \textbf{j}, often lost from the surviving form. \\
Palatal diphthongization & A stressed vowel affected by an initial palatal consonant, especially in West Saxon. \\
Back mutation & A stressed front vowel affected by a following back vowel, with dialect-specific intervening-consonant restrictions. \\
Smoothing & A diphthong becoming a simple vowel in a specified period and dialect. \\
Gemination & A consonant becoming longer; an older \textbf{j} is one possible cause. \\
Vowel loss and contraction & Endings or internal vowels disappearing, or adjacent vowels merging after consonant loss. \\
Ablaut & Inherited stem-vowel alternation, subsequently affected by other changes and analogy. \\
\end{longtable}
\endgroup
% appendix-block:B086:end
% appendix-block:B087:paragraph:begin
The account follows Cook’s selected 1895 second edition of \emph{A First Book in Old English} and R. D. Fulk’s \emph{An Introductory Grammar of Old English with an Anthology of Readings}, in its author-corrected 2020 open-access reprint of the 2014 edition. Fulk’s editorial dots over palatal consonants are omitted in the course-style examples; the pronunciation distinctions are explained in words. The starred earlier forms and Gothic comparison remain explicitly separate from the Old English examples. PDF references below count pages from the beginning of each file, starting at one.\par
% appendix-block:B087:end
% appendix-block:B088:paragraph:begin
1. Fronting and restoration: Fulk §29, printed pp. 18–19 / PDF 30–31; Appendix A IV.26, printed p. 111 / PDF 123. Cook §43.2, printed p. 28 / PDF 44, gives the familiar noun alternation. The back-vowel restoration condition does not extend through arbitrary consonant clusters.\par
% appendix-block:B088:end
% appendix-block:B089:paragraph:begin
2. Breaking: Fulk §62, printed pp. 41–42 / PDF 53–54; §91, printed p. 61 / PDF 73; Appendix A IV.28, printed p. 111 / PDF 123. Cook §21, printed pp. 16–17 / PDF 32–33, explicitly qualifies his shorthand by identifying \textbf{æ} as the immediate input where he first writes \textbf{a}. The vowel-specific restrictions before \textbf{l} follow Fulk. His Appendix A V.45, printed p. 113 / PDF 125, explicitly excludes breaking in unstressed syllables.\par
% appendix-block:B089:end
% appendix-block:B090:paragraph:begin
3. Front mutation: Cook §§15–17, printed pp. 12–13 / PDF 28–29, and Appendix V, printed pp. 268–269 / PDF 284–285; Fulk §74, printed pp. 48–49 / PDF 60–61, §79, printed p. 52 / PDF 64, §82, printed p. 53 / PDF 65, and §109, printed p. 78 / PDF 90. The comparison forms and the selected \textbf{hǣtt} are those of Units 3 and 5.\par
% appendix-block:B090:end
% appendix-block:B091:paragraph:begin
4. Older raising: Cook §17, printed p. 13 / PDF 29, and §104, printed pp. 61–62 / PDF 77–78; Fulk §§73–76, printed pp. 48–49 / PDF 60–61, and Appendix A IV.18, printed p. 110 / PDF 122. Both authors qualify the practical grouping of \textbf{e} beside \textbf{i} with later umlaut.\par
% appendix-block:B091:end
% appendix-block:B092:paragraph:begin
5. Palatal diphthongization: Cook §18, printed pp. 13–14 / PDF 29–30; Fulk §71, printed pp. 45–46 / PDF 57–58, and Appendix B I.5, printed p. 119 / PDF 131. The complete principal-part series and the noun forms agree with Units 5 and 11.\par
% appendix-block:B092:end
% appendix-block:B093:paragraph:begin
6. Back mutation: Fulk §119, printed p. 85 / PDF 97, and Appendix B I.3, printed p. 118 / PDF 130; compare Cook §20, printed pp. 15–16 / PDF 31–32. Fulk supplies the West Saxon restrictions and discusses analogical restoration of a simple vowel. The two spellings compared here occur at separate points in his paragraph.\par
% appendix-block:B093:end
% appendix-block:B094:paragraph:begin
7. Smoothing: Fulk §111d, printed p. 79 / PDF 91; Appendix A IV.35, printed p. 112 / PDF 124; Appendix B III.12, printed p. 119 / PDF 131. Late West Saxon and Anglian smoothing have separate conditions. Cook §19, printed pp. 14–15 / PDF 30–31, is useful for other vowel-spelling correspondences, but is not the authority for this smoothing account.\par
% appendix-block:B094:end
% appendix-block:B095:paragraph:begin
8. Palatal consonants and chronology: Fulk §135a–d, printed pp. 103–104 / PDF 115–116; §71, printed p. 45 / PDF 57. The pronunciation guide, §4, supplies the practical hard-consonant reading of \textbf{cēlan}. Fulk notes disagreement over some of the detailed palatalization conditions.\par
% appendix-block:B095:end
% appendix-block:B096:paragraph:begin
9. Gemination: Fulk §67, printed p. 44 / PDF 56, and Appendix A III.4–6, printed p. 109 / PDF 121; Cook §§33–36, printed pp. 24–25 / PDF 40–41. Cook Appendix V also cautions against treating Gothic as the parent of the other Germanic languages.\par
% appendix-block:B096:end
% appendix-block:B097:paragraph:begin
10. Vowel loss and weight: Cook §4, printed p. 6 / PDF 22; §23, printed pp. 18–19 / PDF 34–35; §43.4, printed pp. 28–29 / PDF 44–45; §§46–47, printed pp. 31–33 / PDF 47–49. Fulk §§32–33, printed p. 21 / PDF 33, and §137a–b, printed p. 104 / PDF 116, supply the historical account and the full-stem variant.\par
% appendix-block:B097:end
% appendix-block:B098:paragraph:begin
11. Loss of \textbf{h} and contraction: Fulk §§131–132, printed pp. 98–99 / PDF 110–111; compare Cook §§27–29, printed pp. 21–23 / PDF 37–39. The imperative is also supplied in Unit 6. Compensatory lengthening and contraction have distinct environments.\par
% appendix-block:B098:end
% appendix-block:B099:paragraph:begin
12. Ablaut: Cook §22, printed pp. 17–18 / PDF 33–34, and §104, printed pp. 61–62 / PDF 77–78; Fulk §85, printed p. 57 / PDF 69, and §91, printed p. 61 / PDF 73. Fulk Appendix A VI offers a fuller account of the inherited vowel system.\par
% appendix-block:B099:end
% appendix-block:B100:paragraph:begin
13. Consonant alternation: Cook §37, printed p. 25 / PDF 41; Fulk §87, printed p. 58 / PDF 70, and Appendix A II, printed pp. 108–109 / PDF 120–121. These locators support a bounded explanation of the historical alternation, rather than a new rule for freely exchanging consonants.\par
% appendix-block:B100:end
\end{document}
